% Generated draft: requires mathematical and source review.
\documentclass[11pt]{article}
\usepackage[margin=25mm]{geometry}
\usepackage{fontspec}
\setmainfont{Latin Modern Roman}
\usepackage{amsmath,amssymb,mathtools,unicode-math}
\setmathfont{Latin Modern Math}
\usepackage{enumitem,needspace,xurl,hyperref}
\setlength{\parindent}{0pt}
\setlength{\parskip}{5pt}
\setlength{\emergencystretch}{3em}
\newcommand{\N}{\mathbb{N}}
\newcommand{\Z}{\mathbb{Z}}
\newcommand{\Q}{\mathbb{Q}}
\newcommand{\R}{\mathbb{R}}
\newcommand{\C}{\mathbb{C}}
\newcommand{\F}{\mathbb{F}}
\newcommand{\A}{\mathbb{A}}
\newcommand{\PP}{\mathbb{P}}
\newcommand{\E}{\mathbb{E}}
\newcommand{\eps}{\varepsilon}
\newcommand{\dd}{\,\mathrm{d}}
\begin{document}
\section*{UnsolvedMath ID: 30003767}
\textbf{Source problem code:} catalogue entry
\subsection{Definitions and mathematical statement}
Let \(f(t,x)\in\Q[t,x]\) be irreducible, as in the supplied record. The parameter domain and the treatment of exceptional specializations are not specified there. The following gives a conditional formulation over rational parameters, with these unresolved choices made explicit.

Suppose \(n=\deg_x f\geq 1\), and write
\[
f(t,x)=\sum_{j=0}^{n}a_j(t)x^j,\qquad a_n(t)\neq 0.
\]
Define the set of parameters preserving the degree in \(x\) by
\[
U_f=\{u\in\Q:a_n(u)\neq 0\}.
\]
For these parameters, define
\[
\operatorname{Red}^{\mathrm{reg}}_f
=\left\{u\in U_f:
\begin{array}{l}
 f(u,x)=g(x)h(x)\text{ for some }g,h\in\Q[x],\\
 \deg g\geq 1\text{ and }\deg h\geq 1
\end{array}\right\}.
\]
This definition concerns reducibility over \(\Q\). It does not determine whether the source's \(\operatorname{Red}_f\) includes parameters at which the degree drops or the specialization is zero or constant. The hypothesis \(\deg_x f\geq1\) and the rational parameter domain are natural conventions for this formulation, but they are not explicit hypotheses of the supplied statement.

One conventional meaning of a thin-set presentation of a set \(S\subseteq\Q\) is a finite covering
\[
S\subseteq E\cup\bigcup_{i=1}^{r}\varphi_i\bigl(C_i(\Q)\bigr),
\]
where \(r\in\Z_{\geq0}\), \(E\subseteq\Q\) is finite, each \(C_i\) is an integral algebraic curve over \(\Q\), and each \(\varphi_i:C_i\to\mathbb A^1_{\Q}\) is a dominant generically finite morphism defined over \(\Q\), with
\[
d_i=\deg\varphi_i=[\Q(C_i):\Q(t)]\geq2.
\]
The source does not specify whether this class of maps is intended, whether the curves must be rational, or whether the displayed containment must instead be equality, possibly outside a finite exceptional set. These alternatives must be resolved before a unique quantitative problem is obtained.

For a specified specialization set \(S\subseteq\Q\), define its symmetric integer counting function by
\[
N_S(B)=\#\{u\in S\cap\Z:|u|\leq B\},\qquad B\in\R,\ B\geq1.
\]
\textbf{Problem, conditional on the conventions above.} For a given irreducible \(f\), determine how many maps \(\varphi_i\) are required to describe its reducible-specialization set, which degree lists \((d_1,\ldots,d_r)\) are required or possible, and how such presentation data constrain the growth of \(N_{\operatorname{Red}_f}(B)\) as \(B\to\infty\). The supplied background also identifies the finiteness of \(\operatorname{Red}_f\cap\Z\) as part of the unresolved inquiry. The source gives no proposed bound, asymptotic formula, criterion for finiteness, or specified dependence on invariants of \(f\). It also does not specify whether the number of maps is to be minimized; arbitrary presentations have no intrinsic number of components.
\subsection{Short English statement}
For an irreducible polynomial over \(\Q\), describe the rational parameters giving reducible specializations by finitely many value maps, and relate the number and degrees of those maps to the growth and possible finiteness of the integer specialization set. The parameter domain, exceptional specializations, permitted maps, and precise meaning of a presentation remain unspecified in the source. This is a statement-only draft requiring mathematical review.
\subsection{Sources}
\begin{itemize}
\item Local UnsolvedMath registry, supplied complete record for permanent ID \(30003767\), titled ``Quantitative Thin-Set Descriptions of Reducible Specializations,'' source problem code \texttt{catalogue entry}, in the Oberwolfach Reports Open Problems collection. No individual proposer is supplied. The record lists proposed year 2018.
\item Supplied report citation: \emph{Mini-Workshop: Arithmetic Geometry and Symmetries around Galois and Fundamental Groups (2019)}, Oberwolfach Reports, pp.~1021--1080, DOI \texttt{10.4171/owr/2018/17}. \url{https://doi.org/10.4171/owr/2018/17}. The supplied literature list calls this the ``OWR 2018 report''; the discrepancy with the citation's 2019 designation is unresolved.
\item Supplied literature reference: Konig--Neftin, ``reducible specializations.'' \url{https://arxiv.org/abs/2001.03630}.
\item Supplied literature reference: Kadets--Neftin, ``Hilbert irreducibility.'' \url{https://arxiv.org/abs/2606.10584}.
\item Supplied literature assessment, checked on \texttt{2026-08-21}: the record labels the problem open, attributes substantial special cases of reducible specialization sets to Konig--Neftin, and says a 2026 paper treats further finite-degree cases while the general number and degrees of value maps and the finiteness problem remain unresolved. This is supplied dated curation derived from the OWR record, not an independently verified literature assessment; its status and cited claims require independent verification.
\item The registry labels the statement \texttt{corrected\_verified} and reports that notation, typography, and omitted definitions were repaired while preserving mathematical content. These are supplied editorial claims, not mathematical verification of this draft. No cited source was visited in preparing this response.
\end{itemize}
\end{document}
